\begin{table}[htbp]
\centering
\caption{Comparison of Formula One race strategy modelling approaches.}
\label{tab:existing_research}

\small
\renewcommand{\arraystretch}{1.25}

\begin{tabularx}{\textwidth}{|
>{\raggedright\arraybackslash}p{0.21\textwidth}|
>{\raggedright\arraybackslash}p{0.18\textwidth}|
>{\raggedright\arraybackslash}p{0.16\textwidth}|
X|
}
\hline
\textbf{Source}
&
\textbf{Method}
&
\textbf{Scope}
&
\textbf{Relevance to this project}
\\
\hline

Bekker \& Lotz
\cite{BekkerLotz2009}
&
Discrete-event simulation
&
Whole-race strategy
&
Demonstrates historical race simulation and validation against real race outcomes, supporting the use of simulation as a decision-support tool.
\\
\hline

Carrasco Heine \& Thraves
\cite{CarrascoHeineThraves2023}
&
Dynamic programming
&
Whole-race optimisation
&
Models tyre wear and cumulative race time. Supports lap time based optimisation but not isolated undercut analysis.
\\
\hline

Piccinotti
\cite{Piccinotti2021}
&
Monte Carlo planning
&
Whole-race planning
&
Uses historical race data and probabilistic simulation to evaluate strategic options under uncertainty.
\\
\hline

Aguad \& Thraves
\cite{AguadThraves2024}
&
Dynamic programming and game theory
&
Two-driver interaction
&
Introduces direct interaction between competing drivers, conceptually similar to the two-driver comparison considered in this work.
\\
\hline

Thomas et al.
\cite{ThomasEtAl2025}
&
Explainable reinforcement learning
&
Whole-race optimisation
&
Demonstrates modern AI-based strategy optimisation while maintaining model interpretability.
\\
\hline

Hettmann
\cite{Hettmann2024}
&
Machine learning
&
Pit stop prediction
&
Predicts pit stop timing from historical races rather than evaluating the engineering causes of successful undercuts.
\\
\hline

Fatima \& Johrendt
\cite{FatimaJohrendt2023}
&
Deep neural networks
&
Race outcome prediction
&
Illustrates the increasing application of deep learning to Formula One strategy, although with limited transparency.
\\
\hline

\end{tabularx}

\end{table}